\section{符号说明}
% 篇幅上限：一页。
% 闭环要求（15Verification 会当硬错误抓）：正文用到的每个符号都在本表里；本表每条都在正文出现过。
%
% 本表三列：符号 / 含义 / 单位，单位单独成列，不塞进「含义」的括号里。
%    （参考件 `基于变物性耦合传热传质与动域模型的.pdf` 即此形制。）
%    写成 `$\rho$ & 湿基密度（kg/m$^3$）` 是错的，必须写成两格：
%    `$\rho$ & 湿基密度 & kg/m$^3$`。同一行有多个量时，单位列用分号并列，与含义列逐项对应：
%    `$c_p,\ k$ & 比热容、热传导系数 & J/(kg$\cdot$K)、W/(m$\cdot$K)`；
%    无量纲写「无量纲」，不写「-」。
%
% **本节不写开头总述段**（不要开头总结）：
%   如「除特别说明外，温度用摄氏度…全文所用符号的含义与单位汇总于表 1。」这种段——
%   那是总述/元话语，不是符号说明本身；读者要的是表。单位口径若确实需要交代，
%   写进表注或正文该量出现处，不在这里铺一段。
%   直接 \section{符号说明} 之后出表（下面那段注释讲的是表体本身怎么排）。

% 本表不作浮动体：它是本节唯一的实体内容，高约 20cm，
% 用 table 浮动环境时会因当页放不下而漂到下一页——会漂到「五、模型的建立与求解」
% 标题之后，看上去像符号表跑进了第 5 节。
% 也不用 center+tabular：那样整表不可分割，而本页顶部常被图占去约 207pt，
% 剩余空间装不下整表，TeX 只能把整张表推到下一页，本节标题与正文之后留下约半页空白。
% 改用 longtable：表体可在行间断页，页面填满且始终留在本节之内。
\setlength{\LTleft}{\fill}
\setlength{\LTright}{\fill}
% 必须显式补一段竖向空白：preamble 为统一行距设了 \lineskip=0pt、\lineskiplimit=-20pt，
% 而 longtable 是直接进竖向列表的盒子，TeX 会在它与上一段最后一行之间插入行间胶，
% 该胶在上述设置下被压成 0pt，表头会直接压在「…汇总于表 1。」那一行上。
% 显式 \vspace 之后上一项变成胶，TeX 不再插行间胶，间距恢复正常。
\par\vspace{0.25\baselineskip}
% 表标题放在 longtable 之外，作为一段居中小字（字号与 \caption 的 capstyle 一致）。
% 放进表内当一行时，该行行高由 \@arstrutbox 决定、又受 \lineskip=0pt 影响，
% 只比下一行高 5.5pt，标题被压在 \midrule 上、与表头重叠。移出表外即可完全避开。
% \nopagebreak 保证标题不会与表体分处两页。
{\centering\fontsize{10.5pt}{12.6pt}\selectfont\bfseries 表 1\quad 符号说明\par}
\nopagebreak
\vspace{0.5em}
{\small
  % 行距补丁：preamble 全局设了 \lineskip=0pt、\lineskiplimit=-20pt（为统一正文行距）。
  % 这套设置在普通 tabular 里没问题，但 longtable 会因此把每一行的行间胶压成 0，
  % 表头行只剩 5.6pt 高、\toprule 横穿「符号 含义 单位」。表内改回正规矩即可；
  % \lineskip 取 \baselineskip，即使回退到它也等于正常行距，不会压行。
  \setlength{\lineskiplimit}{0pt}%
  \setlength{\lineskip}{\baselineskip}%
  % 表头行由 5.6pt 恢复到正常高度后整表高了约 20pt，本表行距单独收到 1.0 以抵消。
  \renewcommand{\arraystretch}{1.0}%
  % 三列对齐：符号靠左、含义居中、单位靠右
  \begin{longtable}{@{}L{0.19\textwidth}C{0.57\textwidth}R{0.19\textwidth}@{}}
    % 表头行不能放进 \endfirsthead：那样 \toprule 会被画到「符号 含义 单位」中间，
    % 且表头下的 \midrule 直接丢失——这是 longtable 首块的手工规则放置问题。
    % 把表头行改成正文首行即可两者兼得：横线位置正确，且续页不重复表头。
    \endfirsthead
    \endhead
    \bottomrule
    \endlastfoot
    \toprule
    符号 & 含义 & 单位 \\
    \midrule
    % ↓↓ 以下两行只是版式示例（照抄会判「符号表与正文对不上」）：全部替换为本题符号。
    %    排列次序照正文出现次序或按类别分组（几何/物性/环境/数值），别按字母表排。
    %    第一行示「一符一义」；第二行示「一格多个量」——含义与单位用顿号/分号逐项对齐。
    %    占位符用 ASCII（`$<S_1>$`）是有意的：汉字进数学模式渲染不出来，留空会看不清版式。
    $<S_1>$ & <含义：只写「是什么」，不写单位> & <单位> \\
    $c_p,\ k$ & 比热容、热传导系数 & J/(kg$\cdot$K)、W/(m$\cdot$K) \\
  \end{longtable}}

% 收尾（可选但常拿分）：把同名不同义的量点清楚，凡有歧义就写这一段。
% 例：「附件 1 的 $C_\infty$ 是空气含湿量，与药材干基含水率 $C$ 基准不同（见模型假设）」
%     「$\bar C_\infty$（平台环境均值）、$C_{lo}$（环境下确界）与 $\bar C$（药材体积平均）三者互不相同」。
